\documentclass[10pt]{article}
\usepackage[T1]{fontenc}
\usepackage[utf8]{inputenc}
\usepackage[french]{babel}
\usepackage{amsmath,amssymb,mathrsfs,mathpazo}
\usepackage{geometry}
\geometry{a4paper,margin=2cm}
\usepackage{graphicx,array,tabularx,tabularray,colortbl,multirow,multicol,booktabs,enumitem,pifont,ulem,fancybox,diagbox,cancel}
\usepackage{pstricks,pstricks-add,pst-plot,pst-tree,pst-node,pst-text,pst-func,pst-math,pst-bspline,pst-3dplot,pst-eucl,pst-solides3d}
\let\pscalculate\undefined
\usepackage{pst-bezier}
\usepackage{tikz,pgfplots,tkz-tab}
\usetikzlibrary{arrows,arrows.meta,patterns,calc,positioning,trees,decorations.markings,shapes.geometric,shapes.misc,shapes.arrows,intersections,decorations.pathmorphing,decorations.pathreplacing,backgrounds,shadows,babel}
\pgfplotsset{compat=1.18}
\usepgfplotslibrary{fillbetween}
\usepackage[np]{numprint}
\usepackage{listings,makecell,pas-tableur,twoopt}
\let\barre\undefined
\usepackage[pstricks]{tablvar}
\definecolor{darkgreen}{RGB}{0,100,0}
\providecommand{\red}{\color{red}}
\providecommand{\blue}{\color{blue}}
\providecommand{\green}{\color{green}}
\providecommand{\black}{\color{black}}
\providecommand{\white}{\color{white}}
\providecommand{\gray}{\color{gray}}
\providecommand{\cyan}{\color{cyan}}
\providecommand{\magenta}{\color{magenta}}
\providecommand{\yellow}{\color{yellow}}
\providecommand{\violet}{\color{violet}}
\providecommand{\orange}{\color{orange}}
\providecommand{\euro}{\texteuro}
\providecommand{\EUR}{\texteuro}
\providecommand{\ieme}{\textsuperscript{e}}
\providecommand{\degres}{\ensuremath{^\circ}}
\providecommand{\degree}{\ensuremath{^\circ}}
\providecommand{\vv}[1]{\overrightarrow{#1}}
\providecommand{\Coord}[1]{#1}
\providecommand{\decosix}{}
\providecommand{\starredbullet}{\ensuremath{\bullet}}
\pagestyle{empty}
\setlength{\parindent}{0pt}
\usepackage{adjustbox}\DeclareUnicodeCharacter{2236}{:}\DeclareUnicodeCharacter{2190}{\ensuremath{\leftarrow}}
\providecommand{\R}{\mathbb{R}}
\providecommand{\N}{\mathbb{N}}
\providecommand{\Z}{\mathbb{Z}}
\providecommand{\Q}{\mathbb{Q}}
\providecommand{\D}{\mathbb{D}}
\providecommand{\C}{\mathbb{C}}
\providecommand{\e}{{\mathrm{e}}}
\providecommand{\vect}[1]{\overrightarrow{#1}}
\providecommand{\vectt}[1]{\overrightarrow{\mathrm{#1}}}
\providecommand{\vvt}[1]{\overrightarrow{\mathrm{#1}}}
\providecommand{\np}[2][]{{#2~#1}}
\providecommand{\vv}[1]{{\overrightarrow{#1}}}
\renewcommand{\d}{{\mathrm{d}}}
\providecommand{\degres}{{^\circ}}
\providecommand{\degre}{{^\circ}}
\providecommand{\blue}{\color{blue}}
\providecommand{\red}{\color{red}}
\providecommand{\enskip}{\quad}
\providecommand{\llbracket}{\lbrack\!\lbrack}
\providecommand{\rrbracket}{\rbrack\!\rbrack}
\providecommand{\euro}{{\text{€}}}
\providecommand{\no}{{\text{n°}}}
\providecommand{\interval}[2]{{[#1;#2]}}
\providecommand{\intervFO}[2]{{[#1;#2[}}
\providecommand{\dotfill}{\ldots}

\providecommand{\ds}{\displaystyle}
\providecommand{\textding}[1]{#1}
\providecommand{\Oij}{\(\mathrm{O};\vec{\imath},\vec{\jmath}\)}
\providecommand{\Oijk}{\(\mathrm{O};\vec{\imath},\vec{\jmath},\vec{k}\)}
\providecommand{\Ouv}{\(\mathrm{O};\vec{u},\vec{v}\)}
\providecommand{\intervFF}[2]{\ensuremath{[#1;#2]}}
\renewcommand{\euro}{\ifmmode\text{\texteuro}\else\texteuro\fi}
\DeclareUnicodeCharacter{00B0}{\ifmmode^{\circ}\else\textdegree\fi}
\DeclareUnicodeCharacter{20AC}{\ifmmode\text{\texteuro}\else\texteuro\fi}
% Énoncés : Éducation nationale. Transcription et corrigés : APMEP, Denis Vergès et contributeurs. Découpage et classement : CoopMaths, Sébastien Lozano, Éric Elter et contributeurs. Adaptation de la présentation : Maths Explorer.
% CC BY-SA : https://coopmaths.fr/www/about/
% https://coopmaths.fr/alea/static/e3c/2020/tex/e3c_2020_00_sujet27_4_cor.tex
\begin{document}
\noindent\textit{APMEP / CoopMaths. CC BY-SA. Présentation adaptée par Maths Explorer.}\par\medskip
\subsection*{1.}
	
	Au départ, il y a 6 jetons rouges sur 10, donc \(p(R_1) = \dfrac{6}{10} = \dfrac{3}{5} = 0,6\).
	
	Si un jeton rouge a été tiré, il reste 5 rouges sur 9 jetons, donc \(p_{R_1}(R_2) = \dfrac{5}{9}\).
	
	Si un jeton noir a été tiré en premier, il reste 6 rouges sur 9 jetons, donc \(p_{\overline{R_1}}(R_2) = \dfrac{6}{9} = \dfrac{2}{3}\).
	
	\begin{center}
		\begin{tikzpicture}
			\node (P_-1_0) at (-2,-1.5) {$\phantom{A}$};
			\node (P_0_0) at (0,-0.5) {$R_1$};
			\draw (P_-1_0) -- (P_0_0) node[midway, above] {$\frac35$};
			\node (P_1_0) at (2,-0) {$R_2$};
			\draw (P_0_0) -- (P_1_0) node[midway, above] {$\frac59$};
			\node (P_1_1) at (2,-1) {$\overline{R_2}$};
			\draw (P_0_0) -- (P_1_1) node[midway, below] {$\frac49$};
			\node (P_0_2) at (0,-2.5) {$\overline{R_1}$};
			\draw (P_-1_0) -- (P_0_2) node[midway, below] {$\frac25$};
			\node (P_1_2) at (2,-2) {$R_2$};
			\draw (P_0_2) -- (P_1_2) node[midway, above] {$\frac23$};
			\node (P_1_3) at (2,-3) {$\overline{R_2}$};
			\draw (P_0_2) -- (P_1_3) node[midway, below] {$\frac13$};
		\end{tikzpicture}
	\end{center}
	
	\subsection*{2.}
	
	\subsubsection*{a.}
	
	On a \(p(A) = p(R_1 \cap R_2) = p(R_1) \times p_{R_1}(R_2) = \dfrac{3}{5} \times \dfrac{5}{9} = \dfrac{3}{9} = \dfrac{1}{3}\).
	
	\subsubsection*{b.}
	
	L'évènement \(\overline{A}\) signifie \(\text{\guillemotleft{} le joueur obtient au plus un jeton rouge \guillemotright{}}\).
	
	\subsubsection*{c.}
	
	On a \(p(\overline{R_1} \cap R_2) = p(\overline{R_1}) \times p_{\overline{R_1}}(R_2) = \dfrac{2}{5} \times \dfrac{2}{3} = \dfrac{4}{15}\).
	
	D'après la loi des probabilités totales :
	\[
	p(R_2) = p(R_1 \cap R_2) + p(\overline{R_1} \cap R_2) = \dfrac{1}{3} + \dfrac{4}{15} = \dfrac{5}{15} + \dfrac{4}{15} = \dfrac{9}{15} = \dfrac{3}{5} = 0,6.
	\]
	
	\subsubsection*{d.}
	
	On a \(p_{\overline{R_2}}(R_1) = \dfrac{p(\overline{R_2} \cap R_1)}{p(R_2)} = \dfrac{p(R_1 \cap \overline{R_2})}{p(R_2)} = \dfrac{\dfrac{3}{5} \times \dfrac{4}{9}}{\dfrac{3}{5}} = \dfrac{4}{9} < \dfrac{4,5}{9} = 0,5\).\\
	 L'affirmation est donc fausse.

\end{document}
